\documentclass[a4paper]{article}
\usepackage[utf8x]{inputenc}    
\usepackage[T1]{fontenc}
\usepackage{graphicx}
\graphicspath{{Figures/}}
\usepackage[frenchb]{babel}
\usepackage[francais,bloc,nowatermark]{automultiplechoice}    
% [nowatermark] pour enlever le texte "PROJET" et le pied de page
\usepackage{multicol}
\usepackage{tkz-tab}
\usepackage{fancyhdr,amssymb,amsmath}
%\DeclareGraphicsRule{.7}{mps}{*}{}

%===================
\begin{document}
%===================

%%% On répartie les quesitons dans différents groupes.
% ---------------------------------
%%% Questions groupe : "catégorie1"
%----------------------------------
\element{categorie1}{
  \begin{question}{Derive1}    
% Question ouverte   
Résoudre les équations et inéquation suivantes:
\begin{multicols}{3}
\begin{enumerate}
\item $2000\times 1,4^x$=$10$
\item $30 x^{21}=500$ 
\item $2^x=100$
\item $x^{4,5}=13$
\item $5^x < 0,0001$ 
\end{enumerate}
\end{multicols}

   \AMCOpen{lines=11,dots=false}{\wrongchoice{NA}\scoring{0}\wrongchoice{AB}\scoring{1}\wrongchoice{B}\scoring{2}\correctchoice{TB}\scoring{3}}
% AB rapporte 1 point
% B rapporte 2 points
% TB rapporte 3 points     
  \end{question}
}

%
\element{categorie2}{
  \begin{question}{Derive2}  
% Question ouverte  
Exprimer les quantités suivantes en fonction de $log(3)$ et/ou $log(2)$.

\begin{multicols}{3}
\begin{enumerate}
\item $log(\dfrac{1}{2})$
\item $log(54)$
\item $log(\dfrac{9}{8})$
\item $log(0,09)$
\item $log(0,0081)$
\item $log(0,3)$
\end{enumerate}
\end{multicols}
\AMCOpen{lines=14,dots=false}{\wrongchoice{NA}\scoring{0}\wrongchoice{AB}\scoring{1}\wrongchoice{B}\scoring{2}\correctchoice{TB}\scoring{3}}
% AB rapporte 1 point
% B rapporte 2 points
% TB rapporte 3 points         
  \end{question}
   
}




\element{categorie4}{
  \begin{question}{variation1}  
% Question ouverte  
Le nombre de bactérie au sein d'un être vivant est modéliser par la fonction $b$ définie sur $\mathbb{R}_+^*$ telle que $b(t)=15000 \times (1,343)^t$ avec $t$ exprimé en heure.

\begin{enumerate}
\item Que cherche t-on à savoir quand on résout l'inéquation $b(t)>50000$ ?
\item Résoudre l'inéquation $b(t)>50000$ et conclure par une phrase faisant le lien avec le problème donné.
\end{enumerate} 


   \AMCOpen{lines=14,dots=false}{\wrongchoice{NA}\scoring{0}\wrongchoice{AB}\scoring{1}\wrongchoice{B}\scoring{2}\correctchoice{TB}\scoring{3}}
% AB rapporte 1 point
% B rapporte 2 points
% TB rapporte 3 points         
  \end{question}
}


\element{categorie4}{
  \begin{question}{variation2}  
% Question ouverte  
Monsieur Truque souhaite se constituer une épargne en plaçant une certaine somme d'argent sur un
compte rémunéré bloqué. Pour savoir quand son objectif sera atteint, il doit résoudre l'inéquation
$$1000 \times (1,035)^n > 1600$$

Déterminer l'objectif de Monsieur Truque, les modalités de son placement puis résoudre le problème.

   \AMCOpen{lines=13,dots=false}{\wrongchoice{NA}\scoring{0}\wrongchoice{AB}\scoring{1}\wrongchoice{B}\scoring{2}\correctchoice{TB}\scoring{3}}
% AB rapporte 1 point
% B rapporte 2 points
% TB rapporte 3 points         
  \end{question}
}

%-------------------------

%%%%%%%%%%%%%%%%%%%%%%%%%
%% fin des questions %%%%
%%%%%%%%%%%%%%%%%%%%%%%%%


%%%%%%%%%%%%%%%%%%%%%%%%%%%%
%%% fabrication des copies%%
%%%%%%%%%%%%%%%%%%%%%%%%%%%%
\exemplaire{1}{    

%%% debut de l'entête des copies :    
	%%% première ligne

	\begin{minipage}{1\linewidth}
	  \centering\large\bf DS de Mathématiques TSTMG sur le logarithme décimal  
	\end{minipage}

	%%% codage des numéros d'étudiants sur 3 chiffres par les candidats et encadré d'écriture manuelle des noms.
%	\vspace*{3mm}
	{
	%\setlength{\parindent}{0pt}\AMCcode{etu}{2}\hspace*{\fill}
	\begin{minipage}[b]{12cm}
	% \hspace*{5mm}
  	%  Codez votre numéro ci contre chiffre par chiffre, puis complétez l'encadré.
    % \vspace{3ex}

	 \hfill\champnom{
	 \centering
	 \fbox{
	 \begin{minipage}{1\linewidth}
		\vspace*{3mm}
		NOM - Prénom - Classe :
		\vspace*{3mm}
	 \end{minipage}
	    }	
	   }
	\hfill\vspace{1ex}
% texte de présentation 
\begin{center}\em

Durée : 55 minutes.

  Aucun document n'est autorisé. La calculatrice est autorisée.\\
  Toutes les réponses doivent être justifiées.
%  Les questions faisant apparaître le symbole \multiSymbole{} peuvent
%  présenter une ou plusieurs bonnes réponses. Les autres ont une unique bonne réponse.
%  Les réponses fausses ou incohérentes retirent des points. Pour éviter les erreurs machines, il est nécessaire de \textbf{noircir vos cases}.

\end{center}

	\end{minipage}\hspace*{\fill}
	}

\vspace{4ex}

%%% Fin de l'entête
%%%%%%%%%%%%%%%%%%%%%%%

%% Composition des QCM
%%% On mélange des questions par groupe
% On efface le {tout} précédent
\cleargroup{tout}

% On mélange par "catégorie", et on prélève aléatoirement [n] questions 
% du lot {categorie-k} vers la copie complète {tout}
% puis on restitue la copie.
% Si on met [0], toutes les questions sont prises.
\melangegroupe{categorie1}\copygroup[0]{categorie1}{tout}
\melangegroupe{categorie2}\copygroup[0]{categorie2}{tout}
%\melangegroupe{categorie3}\copygroup[0]{categorie3}{tout}
\melangegroupe{categorie4}\copygroup[0]{categorie4}{tout}

\restituegroupe{tout}

% Saut d'une page si le sujet comporte un nombre impair de pages.
% Ainsi il est possible d'imprimer les sujets en recto-verso.

\AMCcleardoublepage    

%%% Fin de la création des questionnaires

}
% Fin de exemplaires{k}{ 

\end{document}

